\documentclass[aspectratio=169]{beamer}
\usetheme{Madrid}
\usecolortheme{default}
\usepackage[T1]{fontenc}
\usepackage[utf8]{inputenc}
\usepackage{lmodern}
\usepackage{booktabs}
\usepackage{array}
\usepackage{amsmath}
\usepackage{graphicx}
\usepackage{ragged2e}
\usepackage{multicol}

\title{StockSense AI}
\subtitle{XGBOOST-BASED FUNDAMENTAL STOCK SCREENING AND LONG-TERM INVESTMENT RECOMMENDATION SYSTEM}
\author{Syed Ahmed}
\institute{Mini Project / Hackathon Presentation}
\date{\today}

\begin{document}

\begin{frame}
  \titlepage
\end{frame}

\begin{frame}{Problem Statement}
\justifying
Long-term investors usually check whether a company is fundamentally strong before buying a stock. This requires reading many financial ratios such as ROE, ROCE, debt-to-equity, P/E, growth, free cash flow, promoter holding, and pledge percentage. For beginners and family investors, this process is slow, confusing, and dependent on personal experience.

\vspace{0.3cm}
\textbf{Project goal}
\begin{itemize}
  \item collect important company fundamentals,
  \item analyze long-term fundamental strength,
  \item generate a 0--100 AI strength score,
  \item classify stocks as Buy, Watchlist, or Avoid,
  \item provide beginner-friendly guidance and risk warnings.
\end{itemize}
\end{frame}

\begin{frame}{What Was Implemented from Research Papers}
\small
\begin{itemize}
  \item From fundamental-analysis ML papers: use financial indicators instead of only price charts.
  \item From stock recommendation papers: convert financial indicators into ranking/recommendation output.
  \item From finance-ML review papers: consider model evaluation, data quality, and interpretability.
  \item From practical decision-support systems: provide a user-facing workflow, not only a model result.
\end{itemize}

\vspace{0.2cm}
\textbf{What is added beyond the papers}
\begin{itemize}
  \item beginner guide videos for collecting values,
  \item 15-parameter web-based input form,
  \item XGBoost-powered score generation,
  \item explainable category-level diagnostics,
  \item macro-risk alerts for wars, tariffs, rate hikes, and policy shocks.
\end{itemize}
\end{frame}

\begin{frame}{Research Paper Comparison}
\small
\begin{tabular}{p{3.4cm} p{4.7cm} p{2.1cm} p{3.0cm}}
\toprule
\textbf{Paper} & \textbf{Core idea} & \textbf{Data type} & \textbf{Relevance to project} \\
\midrule
Huang et al. (IEEE SSCI 2021) & ML for stock prediction based on fundamental analysis & Quarterly financial data & Supports use of fundamentals for ML-assisted investment decisions \\
\addlinespace
Yang et al. (arXiv 2025) & Dynamic stock recommendation using financial indicators and model selection & Stock indicators and returns & Matches recommendation/ranking direction \\
\addlinespace
Saberironaghi et al. (AppliedMath 2025) & Review of ML and deep learning in stock market prediction & Survey paper & Helps justify ML use and discuss limitations \\
\bottomrule
\end{tabular}

\vspace{0.25cm}
\textbf{Project positioning:} StockSense AI is closer to fundamental stock screening and recommendation than short-term stock price prediction.
\end{frame}

\begin{frame}{Dataset Used}
\small
\justifying
The project currently uses a synthetic expert-rule dataset designed to imitate multiple long-term stock fundamental profiles. Each row represents one company-year style snapshot with 15 financial indicators and one target score.

\vspace{0.2cm}
\textbf{Processed dataset summary}
\begin{itemize}
  \item total rows: \textbf{62,500}
  \item feature columns: \textbf{15}
  \item target column: \textbf{score} from 0 to 100
  \item output classes: \textbf{Buy, Watchlist, Avoid}
\end{itemize}

\textbf{Score distribution}
\begin{itemize}
  \item Avoid \((<60)\): \textbf{24,937}
  \item Watchlist \((60--79)\): \textbf{18,820}
  \item Buy \((\geq80)\): \textbf{18,743}
\end{itemize}
\end{frame}

\begin{frame}{Synthetic Dataset Coverage}
\small
\justifying
The dataset was strengthened to handle many real-world-like cases, not only ideal examples.

\vspace{0.2cm}
\begin{multicols}{2}
\begin{itemize}
  \item quality compounders,
  \item average companies,
  \item expensive quality stocks,
  \item cyclicals,
  \item distressed companies,
  \item debt traps,
  \item pledge-risk companies,
  \item turnaround cases,
  \item negative cash-flow growth cases,
  \item zero-debt high coverage cases.
\end{itemize}
\end{multicols}

\vspace{0.2cm}
\textbf{Important note:} The dataset is strong for hackathon demonstration and repeatable training, but real-world investment-performance claims require historical fundamentals and future return validation.
\end{frame}

\begin{frame}{Input Parameters}
\small
\begin{columns}
\column{0.33\textwidth}
\textbf{Valuation}
\begin{itemize}
  \item Stock P/E
  \item Industry P/E
  \item Price to book value
\end{itemize}

\textbf{Profitability}
\begin{itemize}
  \item ROE
  \item ROCE
  \item OPM
\end{itemize}

\column{0.33\textwidth}
\textbf{Safety}
\begin{itemize}
  \item Debt to equity
  \item Interest coverage
  \item Current ratio
\end{itemize}

\textbf{Growth}
\begin{itemize}
  \item Sales growth 5 years
  \item Profit growth 5 years
  \item EPS growth 3 years
\end{itemize}

\column{0.33\textwidth}
\textbf{Cash and ownership}
\begin{itemize}
  \item Free cash flow
  \item Promoter holding
  \item Promoter pledge
\end{itemize}
\end{columns}
\end{frame}

\begin{frame}{End-to-End Architecture}
\small
\begin{center}
\begin{tabular}{p{3.1cm} p{7.7cm}}
\toprule
\textbf{Layer} & \textbf{Function} \\
\midrule
Frontend & Guide videos, 15-input analysis form, progress indicator, result gauge, recommendation page \\
\addlinespace
Flask backend & Validates inputs, loads feature order, prepares model input, handles errors, exposes health route \\
\addlinespace
ML engine & XGBoost regressor predicts a 0--100 fundamental strength score \\
\addlinespace
Explanation layer & Groups signals into valuation, profitability, balance-sheet safety, growth, and ownership/cash quality \\
\addlinespace
Risk guidance & Warns users to re-evaluate during wars, tariffs, interest-rate hikes, economic shocks, or major policy changes \\
\bottomrule
\end{tabular}
\end{center}

\vspace{0.25cm}
\centering
User fundamentals $\rightarrow$ Flask validation $\rightarrow$ XGBoost model $\rightarrow$ Score $\rightarrow$ Buy / Watchlist / Avoid
\end{frame}

\begin{frame}[shrink=8]{Preprocessing, Training and Evaluation Pipeline}
\small
\textbf{Dataset generation}
\begin{enumerate}
  \item define fundamental indicators and expert-rule relations,
  \item generate multiple company regimes and edge cases,
  \item assign a 0--100 score using weighted rules and red-flag penalties,
  \item store final dataset in \texttt{data/training\_data.csv}.
\end{enumerate}

\textbf{Training}
\begin{enumerate}
  \item load 15 feature columns,
  \item clean missing and infinite values,
  \item split into train and test data,
  \item train XGBoost regressor,
  \item save \texttt{stocksense\_model.pkl}.
\end{enumerate}

\textbf{Evaluation}
\begin{enumerate}
  \item compute MAE and \(R^2\),
  \item convert score into Buy, Watchlist, Avoid buckets,
  \item compute bucket accuracy,
  \item export model metadata and feature importance.
\end{enumerate}
\end{frame}

\begin{frame}[shrink=8]{Scoring Logic}
\small
\justifying
The target score is generated using a weighted fundamental-analysis philosophy. Positive signals increase the score, while red flags reduce the score.

\vspace{0.2cm}
\textbf{Positive signals}
\begin{itemize}
  \item high ROE and ROCE,
  \item consistent sales, profit, and EPS growth,
  \item low debt-to-equity,
  \item strong interest coverage,
  \item positive free cash flow,
  \item high promoter holding and low promoter pledge.
\end{itemize}

\textbf{Red flags}
\begin{itemize}
  \item excessive valuation compared with industry P/E,
  \item high leverage,
  \item weak interest coverage,
  \item negative free cash flow,
  \item high promoter pledge,
  \item poor liquidity or weak profitability.
\end{itemize}
\end{frame}

\begin{frame}[shrink=8]{XGBoost Model}
\small
\textbf{Model used:} \texttt{XGBRegressor}

\vspace{0.2cm}
\textbf{Why XGBoost was selected}
\begin{itemize}
  \item works well on structured tabular financial data,
  \item captures non-linear relationships between fundamentals,
  \item handles interactions such as high growth with high valuation,
  \item gives feature importance for explainability,
  \item fast enough for web application inference.
\end{itemize}

\vspace{0.2cm}
\textbf{Training configuration}
\begin{itemize}
  \item estimators: 650
  \item max depth: 4
  \item learning rate: 0.04
  \item subsample: 0.85
  \item column sample by tree: 0.85
  \item regularization: alpha 0.05, lambda 1.5
\end{itemize}
\end{frame}

\begin{frame}{How Recommendation is Generated}
\small
\justifying
The model predicts a continuous score from 0 to 100. The web application then converts the score into an investor-friendly recommendation.

\vspace{0.3cm}
\begin{center}
\begin{tabular}{c c p{6.4cm}}
\toprule
\textbf{Score range} & \textbf{Output} & \textbf{Meaning} \\
\midrule
\(\geq 80\) & Buy & Strong long-term candidate according to entered fundamentals \\
\addlinespace
60--79 & Watchlist & Moderate fundamentals; monitor and re-analyze later \\
\addlinespace
\(<60\) & Avoid & Weak or risky fundamentals; avoid unless fundamentals improve \\
\bottomrule
\end{tabular}
\end{center}
\end{frame}

\begin{frame}{Current XGBoost Analytics}
\small
\textbf{From \texttt{model\_metadata.json}}
\begin{itemize}
  \item Rows after cleaning: \textbf{62,500}
  \item Model: \textbf{XGBRegressor}
  \item Mean Absolute Error: \textbf{2.80}
  \item \(R^2\) score: \textbf{0.982}
  \item Buy/Watchlist/Avoid bucket accuracy: \textbf{91.45\%}
\end{itemize}

\vspace{0.2cm}
\textbf{Interpretation}
\begin{itemize}
  \item MAE of 2.80 means the predicted score is usually within about 3 points of the synthetic label.
  \item \(R^2=0.982\) indicates strong fit to the generated fundamental scoring logic.
  \item Bucket accuracy measures whether the recommendation class is correct.
\end{itemize}
\end{frame}

\begin{frame}{Feature Importance}
\small
\begin{tabular}{p{4.2cm} c p{6.0cm}}
\toprule
\textbf{Feature} & \textbf{Importance} & \textbf{Interpretation} \\
\midrule
ROCE & 0.302 & Most important capital-efficiency signal \\
ROE & 0.211 & Strong shareholder-return signal \\
Profit growth 5Y & 0.115 & Measures long-term business growth \\
Interest coverage & 0.076 & Captures debt servicing safety \\
Sales growth 5Y & 0.073 & Revenue consistency signal \\
Current ratio & 0.051 & Liquidity and short-term safety \\
Free cash flow & 0.044 & Cash-generation quality \\
\bottomrule
\end{tabular}

\vspace{0.25cm}
\textbf{Interpretation:} The model focuses mainly on profitability, capital efficiency, growth, and balance-sheet safety, which matches long-term fundamental analysis.
\end{frame}

\begin{frame}{Website Flow}
\small
\begin{enumerate}
  \item User opens StockSense AI homepage.
  \item User watches guide videos to understand how to collect values.
  \item User enters all 15 fundamental parameters.
  \item Flask validates every value.
  \item XGBoost model generates a score.
  \item Result page shows:
  \begin{itemize}
    \item score gauge,
    \item Buy / Watchlist / Avoid recommendation,
    \item strength signals,
    \item review-before-buying signals,
    \item risk alerts and news-monitoring guidance.
  \end{itemize}
\end{enumerate}
\end{frame}

\begin{frame}{How This Project Improves on Basic ML Demos}
\small
\begin{itemize}
  \item Many stock ML projects focus only on price prediction.
  \item This project focuses on fundamental long-term screening.
  \item It combines:
  \begin{itemize}
    \item financial-ratio based input,
    \item XGBoost model prediction,
    \item beginner-friendly frontend,
    \item explanation layer,
    \item model metadata and feature importance,
    \item risk warnings and re-analysis guidance.
  \end{itemize}
  \item So the project is an end-to-end decision-support application, not only a model notebook.
\end{itemize}
\end{frame}

\begin{frame}{Important Limitation}
\small
\begin{alertblock}{Do not overclaim}
The current model is trained on a synthetic expert-rule dataset. Therefore, the accuracy shows how well the model learned the generated fundamental scoring logic. It does not prove real-world stock return prediction accuracy.
\end{alertblock}

\vspace{0.2cm}
\textbf{Correct claim}
\begin{itemize}
  \item The XGBoost model achieved \textbf{91.45\%} Buy/Watchlist/Avoid bucket accuracy on a 62,500-row synthetic fundamental-analysis dataset.
\end{itemize}

\textbf{Incorrect claim}
\begin{itemize}
  \item The model predicts real stock market returns with 91.45\% accuracy.
\end{itemize}
\end{frame}

\begin{frame}{Future Improvements}
\begin{itemize}
  \item collect real historical fundamentals for NSE-listed companies,
  \item add 1-year, 3-year, and 5-year future return labels,
  \item validate the model on unseen real companies,
  \item add sector-wise normalization,
  \item add peer comparison,
  \item add SHAP-based model explanations,
  \item integrate compliant finance APIs where available,
  \item improve UI with portfolio watchlist and saved analysis history.
\end{itemize}
\end{frame}

\begin{frame}{Conclusion}
\justifying
StockSense AI is a fundamental-analysis based stock screening system for long-term investors. It converts 15 important financial indicators into a 0--100 strength score using an XGBoost model and presents the output as Buy, Watchlist, or Avoid.

\vspace{0.3cm}
The project includes a beginner-friendly Flask web application, guide videos, input validation, explainable category diagnostics, feature-importance reporting, and macro-risk warnings. The current XGBoost model performs strongly on a 62,500-row synthetic expert-rule dataset.

\vspace{0.3cm}
The next major improvement is to validate the same system on real historical fundamentals and future stock returns, which would make the project stronger for research and real-world investment decision support.
\end{frame}

\begin{frame}{References}
\footnotesize
\begin{itemize}
  \item Y. Huang, L. F. Capretz and D. Ho, ``Machine Learning for Stock Prediction Based on Fundamental Analysis,'' IEEE SSCI, 2021.
  \item H. Yang, X.-Y. Liu and Q. Wu, ``A Practical Machine Learning Approach for Dynamic Stock Recommendation,'' arXiv:2511.12129.
  \item M. Saberironaghi, J. Ren and A. Saberironaghi, ``Stock Market Prediction Using Machine Learning and Deep Learning Techniques: A Review,'' AppliedMath, 2025.
\end{itemize}
\end{frame}

\begin{frame}
  \centering
  \Huge Thank You
\end{frame}

\end{document}
